\section{Related Work}
\label{sec:related-work}

\subsection{Algorithmic Stability and Generalization}

The modern theory of algorithmic stability was formalized by Bousquet and
Elisseeff \cite{bousquet2002stability}, who defined several stability notions
centered on replace-one perturbations and fresh-test evaluation, and used them
to derive generalization bounds. Kutin and Niyogi \cite{kutin2002almost}
expanded the landscape with weaker and almost-everywhere variants, while
Feldman and Vondr\'ak \cite{feldman2019high} and later Bousquet, Klochkov, and
Zhivotovskiy \cite{bousquet2020sharpening} obtained sharp high-probability
bounds.  Subsequent work connected stability to learnability and uniform
convergence, clarifying when stability is not merely sufficient but also
structurally related to learnability \cite{shalev2010learnability}. Mukherjee,
Niyogi, Poggio, and Rifkin \cite{mukherjee2006learning} further established that
stability is sufficient for generalization and, for empirical risk minimization,
necessary and sufficient for consistency. Hardt,
Recht, and Singer \cite{hardt2016train} further showed that algorithmic
stability explains generalization for optimization-based learning rules such as
stochastic gradient descent. Our work is orthogonal to this line: we do not
derive generalization bounds, but instead provide a microscopic perturbation
calculus that clarifies which stability notion is being discussed in the first
place.  The classical literature mainly uses replace-one or delete-one with
fresh-test evaluation; we show that LOO evaluation (common in practice) can
give a qualitatively different picture.

\subsection{Nearest-Neighbor Classification and Deleted Estimates}

The nearest-neighbor line begins with Fix and Hodges \cite{fix1951discriminatory}
and the classical risk and consistency results of Cover and Hart
\cite{cover1967nearest,cover1968estimation}, together with the consistency
framework formalized by Stone \cite{stone1977consistent}.  Those works study
statistical performance under distributional sampling rather than worst-case
single-perturbation effects.  Closer to our setting are the deleted-estimate
and local-rule analyses of Rogers and Wagner \cite{rogers1978finite} and
Devroye and Wagner \cite{devroye1979distribution,devroye1979local}, which
establish finite-sample bounds for local discrimination rules using
leave-one-out style evaluation.  These are important precedents, but their
goal is to bound expected error, not to separate LOO from adversarial
replace-one at the level of fixed finite metric witnesses.

\subsection{Leave-One-Out and Conformal Stability}

Kearns and Ron \cite{kearns1999algorithmic} connected algorithmic stability
with leave-one-out cross-validation, already emphasizing that CV-style
perturbation has its own logic.  Recent conformal prediction work makes related
distinctions operational, often in service of coverage certificates rather than
0--1 prediction loss.  Examples include training-conditional guarantees
\cite{liang2023algorithmic,pournaderi2024training} and related stability
analyses.  Our results complement this line by showing
that even the simpler 0--1 loss setting already exhibits a structural gap
between LOO and replacement-based control, before any distributional or
coverage considerations enter.

\subsection{Finite Metric Witnesses and Local Learning Rules}

Metric-space analyses of $k$-NN typically appear in consistency or geometric
generality results \cite{collins2020consistency} rather than
perturbation-separation examples.  To our knowledge, no prior work exhibits
a finite graph-metric sample whose explicit purpose is to show that
fixed-sample LOO stability need not control fixed-sample adversarial replace-one
stability for deterministic $k$-NN under the conventions used here.  The
separation we prove is new in its specificity.

\subsection{What the Paper Achieves}

The paper's distinctive contribution lies in combining (a) a systematic
perturbation calculus that separates perturbation operations from evaluation
points, (b) an exact uniform calculus relating those operations, (c) a general
separation theorem valid for all odd $k \ge 3$ (and even $k$ under
tie-breaking), (d) an abstract framework extending beyond $k$-NN, and
(e) reproducible experimental evidence that the separation occurs on real and
synthetic data.  Earlier computational search infrastructure
\cite{codex2025codex} provided the $1$-NN witness and candidate patterns for
larger $k$, but left the general odd-$k$ statement as a conjecture.  The present
paper upgrades that conjecture to a theorem.
